\documentclass[11pt]{article}
\usepackage[margin=1in]{geometry}
\usepackage{amsmath,amssymb,amsthm}
\usepackage{booktabs}
\usepackage{hyperref}
\newtheorem{theorem}{Theorem}
\newtheorem{corollary}[theorem]{Corollary}
\newtheorem{lemma}[theorem]{Lemma}
\theoremstyle{remark}
\newtheorem{remark}[theorem]{Remark}

\newcommand{\Lind}{\mathcal{L}}
\newcommand{\dn}[1]{\lVert #1 \rVert_\diamond}
\newcommand{\tn}[1]{\lVert #1 \rVert_1}
\newcommand{\ket}[1]{\lvert #1 \rangle}
\newcommand{\bra}[1]{\langle #1 \rvert}
\newcommand{\proj}[1]{\lvert #1 \rangle\langle #1 \rvert}
\newcommand{\Sym}{\mathrm{Sym}}

\title{The dimension factor in sample-based Lindbladian simulation\\is necessary for every protocol}
\author{Sebastian Senge\\Independent Researcher}
\date{26 November 2026}

\begin{document}
\maketitle

\begin{abstract}
In sample-based Lindbladian simulation, a jump operator $L$ on a $d$-dimensional system is given through copies of a program state that encodes $L$, and the task is to apply $e^{t\Lind_L}$ to error $\varepsilon$ in diamond distance. The best known protocol uses $O(d\,t^2/\varepsilon)$ copies, while the best known lower bound for all protocols, $\Omega(t^2/\varepsilon)$, does not depend on $d$. We show that the factor $d$ cannot be avoided.\\
For $d \ge 3$, every protocol, with any number of copies processed jointly, needs at least $(d-2)(1-e^{-t})^2/(32\varepsilon)$ copies whenever $\varepsilon \le (1-e^{-t})/8$; for $t \le 1$ and $\varepsilon \le 0.079\,t$ this is at least $0.0124\,(d-2)\,t^2/\varepsilon$.\\
For $d \ge 2$ and high accuracy, $\varepsilon \le t^2/(C d\log(2d/t))$ and $t \le c$, every protocol needs at least $\tfrac{d-1}{4}(1-O(t))\,t^2/\varepsilon$ copies, and for $t \le c/d$ this holds at every accuracy up to a factor $1-O(dt)$.\\
For wave matrix Lindbladization we prove the upper bound $(d/2 + e^{1/d})\,t^2/\varepsilon$ and a worst-case lower bound $\tfrac{d-1}{4}\,t^2/\varepsilon$.\\
The first bound comes from a dual certificate for universal conjugation: the program of $L = \ket{0}\bra{\psi}$ carries $\bar\psi$, while the target needs $\psi$. The second combines positivity of the Choi matrix with Bernstein's inequality for trigonometric polynomials and a decomposition of the program states into spin components. The scalar steps of the first bound and of the upper bound are checked in Lean.
\end{abstract}

\section{Introduction}

In sample-based simulation the generator of the dynamics is given as quantum data, not as a circuit. For Hamiltonians, density matrix exponentiation applies $e^{-i\rho t}$ from copies of $\rho$, with a number of copies of order $t^2/\varepsilon$ [2]. For a Lindbladian with one jump operator $L$, wave matrix Lindbladization (WML) encodes $L$ into the program state $\ket{L} = (L \otimes I)\ket{\Gamma}$, $\ket{\Gamma} = \sum_j \ket{j}\ket{j}$, with $\lVert L\rVert_2 = 1$, and applies $e^{t\Lind_L}$, $\Lind_L(\rho) = L\rho L^\dagger - \tfrac12\{L^\dagger L, \rho\}$, by $n$ short interactions with fresh copies (see [4] and the work it extends). Its sample complexity was shown to be $O(d^2 t^2/\varepsilon)$ [2] and then $O(d\,t^2/\varepsilon)$, with leading constant $(2d+3)/8$ [1]. An explicit rank-one operator shows that WML itself needs $\Omega(d\,t^2/\varepsilon)$ copies [1, Theorem 10]. The lower bound for all protocols, however, is $\Omega(t^2/\varepsilon)$ with no dependence on $d$, and the general analysis is restricted to $d = 2$ [2, Theorem 5]. In [1, Sec.~VII] it is conjectured that the dimension dependence cannot be removed by any protocol.

We show that the factor $d$ is necessary for every protocol, for $d \ge 3$ up to the constant, and with the constant $(d-1)/4$ in the high-accuracy regime. A protocol with $n$ copies is any channel $\Lambda$ from $S \otimes P^{\otimes n}$ to $S$, where $S = \mathbb{C}^d$ is the system and $P = \mathbb{C}^d \otimes \mathbb{C}^d$ holds one program state $\pi_L = \proj{L}$. Its error on $L$ is $\tfrac12 \dn{\Lambda(\cdot \otimes \pi_L^{\otimes n}) - e^{t\Lind_L}}$.

\begin{theorem}[conjugation bound]\label{thm:conj}
Let $d \ge 3$, $n \ge 1$ and $t > 0$. Every channel with error at most $\varepsilon \le (1-e^{-t})/8$ on all $L = \ket{0}\bra{\psi}$ uses
\[
n \ \ge\ \frac{(d-2)(1-e^{-t})^2}{32\,\varepsilon}
\]
copies. In particular $n \ge 0.0124\,(d-2)\,t^2/\varepsilon$ for $t \le 1$ and $\varepsilon \le 0.079\,t$, and $n \ge 0.0124\,(d-2)/\varepsilon$ for $t \ge 1$ and $\varepsilon \le 0.079$.
\end{theorem}

\begin{theorem}[sharp constant at high accuracy]\label{thm:sharp}
There are absolute constants $c, C$ such that for $d \ge 2$ and $t \le c$, every channel with error at most $\varepsilon \le t^2/(C d \log(2d/t))$ on all $L = \ket{0}\bra{\psi}$ uses at least $\tfrac{d-1}{4}(1 - Ct)\,t^2/\varepsilon$ copies. For $t \le c/d$ the same holds for every $\varepsilon$ with $1 - Ct$ replaced by $1 - Cdt$.
\end{theorem}

Theorem~\ref{thm:conj} covers moderate accuracy and every time, with a small constant and the factor $d-2$. Theorem~\ref{thm:sharp} gives the constant $(d-1)/4$ at high accuracy or short times, for every $d \ge 2$; this is also the error constant of WML on the operator $\proj{0}$ (Theorem~\ref{thm:wmllower}). Neither result assumes any structure of the protocol: the $n$ copies may be processed jointly, with any ancilla.

The mechanism of Theorem~\ref{thm:conj} is conjugation. For $L = \ket{0}\bra{\psi}$ the program is $\ket{0} \otimes \ket{\bar\psi}$, while the coherence that $e^{t\Lind_L}$ creates on the input $\proj{0}$ points along $\psi$. A protocol therefore has to conjugate the $d-1$ amplitudes of $\psi$ orthogonal to $\ket{0}$, and a twirled dual certificate turns the cost of universal conjugation into the bound. Theorem~\ref{thm:sharp} uses a different argument. On the Choi matrix of $\Lambda$, a $2\times 2$ positivity test captures the population that leaves $\ket{0}$, and Bernstein's inequality bounds the derivatives that enter it by the worst-case error. The degree of the trigonometric polynomials involved is reduced from $2n$ to about $\sqrt{n\log(n/t)}$ by a decomposition of $\pi_v^{\otimes n}$ into spin components.

For WML itself we give the per-step bound $\Delta^2\lVert L\rVert_\infty^2 (d/2 + e^{1/d}\lVert L\rVert_\infty^2)$, so $(d/2 + e^{1/d})\,t^2/\varepsilon$ copies suffice, and an operator whose one-step error exceeds the bound of [1, Lemma 4] for $d \ge 5$ (Remark~\ref{rem:gap}). The worst-case leading constant of WML lies between $d/2 - 3/4$ and $d/2 + 1$.

\paragraph{Formal verification.} The scalar steps of Theorems~\ref{thm:conj} and~\ref{thm:wmlupper} are proved in Lean 4 (toolchain v4.36.0-rc1, Mathlib revision \texttt{d875a8de}), with only the standard axioms, in the folder \path{laif/wml-sample-constants/lean} of the repository \url{https://github.com/ssenge/LAIF}. These are the dimension and binomial identities, positivity of the $2\times 2$ blocks, the trace of the certificate, the optimization over the multiplier and the constants of the corollaries; for Theorem~\ref{thm:wmlupper} the properties of $\varphi$, the coefficient bounds, the tail inequality and their sum. The block decomposition of the certificate from the Haar moments, the coherence bound in the proof of Theorem~\ref{thm:conjtheta}, the diamond-norm estimates and all of Theorem~\ref{thm:sharp} are proved on paper only. Scripts with asserts in the same folder check the remaining steps numerically for small $d$ and $n$.

\section{The conjugation bound}\label{sec:conj}

We prove a form of Theorem~\ref{thm:conj} with a free angle. Write $c = \cos\theta$, $s = \sin\theta$, $a = 1 - e^{-t/2}$ and $\gamma = a(1 - ac^2)$. For a unit vector $u \perp \ket{0}$ let $\psi_u = c\ket{0} + s\bar u$, where $\bar u$ is the entrywise conjugate in the computational basis, and let $\varepsilon_\theta$ be the worst case over $u$ of the error of $\Lambda$ on $L_u = \ket{0}\bra{\psi_u}$.

\begin{theorem}\label{thm:conjtheta}
For $d \ge 3$, $n \ge 1$, $t > 0$, $\theta \in (0, \pi/2)$ and every channel $\Lambda$,
\[
n\,\varepsilon_\theta \ \ge\ (d-2)\,\bigl(c\gamma - \varepsilon_\theta/s\bigr)_+^2 .
\]
The bound holds already for the trace-norm error on the single input $\proj{0}$.
\end{theorem}

\begin{proof}
\emph{Reduction.} Fix the input $\proj{0}$. Then $T = \Lambda(\proj{0} \otimes \cdot)$ is a channel on $(\mathbb{C}^d)^{\otimes n}$, here identified with the second factors of the programs, with $\tfrac12\tn{T(\phi_u^{\otimes n}) - \sigma_u} \le \varepsilon_\theta$, where $\sigma_u = e^{t\Lind_{L_u}}(\proj{0})$ and the program is $\ket{L_u} = \sum_j L_u\ket{j}\ket{j} = \ket{0} \otimes \ket{\phi_u}$ with $\phi_u = c\ket{0} + s u$. The program carries $u$; the target carries $\bar u$.

\emph{The target coherence.} The state $\sigma_u$ is supported on $\mathrm{span}\{\ket{0}, \bar u\}$ and its coherence $q = \bra{\bar u}\sigma_u\ket{0}$ is real. Every jump resets the state to $\ket{0}$, and the evolution without jumps maps $\ket{0}$ to $(1-ac^2)\ket{0} - acs\,\ket{\bar u}$. Hence $q = q_0 + c^2\int_0^t e^{-\tau} q(t-\tau)\,d\tau$ with $q_0 = -cs\gamma$, where $c^2e^{-\tau}$ is the density of the first jump. Iterating gives $q \le q_0 < 0$, so $\lvert q\rvert \ge cs\gamma$.

\emph{Two tests.} Let $W_u = -(\ket{0}\bra{\bar u} + \ket{\bar u}\bra{0})$ and $\Pi_u = I - \proj{0} - \proj{\bar u}$, let $u$ be Haar distributed on the unit sphere of $\ket{0}^\perp$, and set $X = \mathbb{E}_u \operatorname{Tr} W_u T(\phi_u^{\otimes n})$ and $Y = \mathbb{E}_u \operatorname{Tr} \Pi_u T(\phi_u^{\otimes n})$. Since $\lVert W_u\rVert_\infty = 1$, $\operatorname{Tr} W_u\sigma_u = -2q$ and $\operatorname{Tr}\Pi_u\sigma_u = 0$, we have $X \ge 2cs\gamma - 2\varepsilon_\theta$ and $Y \le \varepsilon_\theta$.

\emph{A dual certificate.} Let $J_T = \sum_{ij} T(\ket{i}\bra{j}) \otimes \ket{i}\bra{j}$, so that $\operatorname{Tr} M\,T(\rho) = \operatorname{Tr} J_T(M \otimes \rho^{\mathsf T})$, $J_T \ge 0$ and $\operatorname{Tr}_{\mathrm{out}} J_T = I$. The transpose maps $\phi_u$ to $\chi_v = c\ket{0} + s v$ with $v = \bar u$, which is again Haar distributed. So for every $\lambda > 0$,
\[
X - \lambda Y = \operatorname{Tr} J_T\,\Omega, \qquad \Omega = \mathbb{E}_v\bigl[(-\ket{0}\bra{v} - \ket{v}\bra{0} - \lambda\Pi_v) \otimes \proj{\chi_v}^{\otimes n}\bigr].
\]
This is where the conjugation enters: both factors of $\Omega$ now carry $v$. Let $H_k \subset (\mathbb{C}^d)^{\otimes n}$ be the span of the normalized symmetric states $\ket{k;v}$ with $k$ factors $v$ and $n-k$ factors $\ket{0}$; then $H_k \cong \Sym^k(\mathbb{C}^{d-1})$, with projector $\Pi_k$ and dimension $D_k = \binom{k+d-2}{k}$. We have $\chi_v^{\otimes n} = \sum_k \sqrt{b_k}\,\ket{k;v}$ with the binomial weights $b_k = \binom{n}{k}s^{2k}c^{2(n-k)}$, and $\mathbb{E}_v (vv^\dagger)^{\otimes k} = \Pi_{\Sym^k}/D_k$. Averaging over the phase $v \mapsto e^{i\alpha}v$ leaves only the cross terms $\ket{0}\bra{v} \otimes \ket{k;v}\bra{k-1;v}$, their adjoints, and terms diagonal in $k$, and $\Omega$ vanishes on $\ket{0}\otimes H_0$ and outside $\bigoplus_k \mathbb{C}^d \otimes H_k$. The cross term maps the copy of $\Sym^k(\mathbb{C}^{d-1})$ inside $\mathbb{C}^{d-1} \otimes H_{k-1}$ isometrically onto $\ket{0} \otimes H_k$ with weight $g_k = \sqrt{b_k b_{k-1}}/D_k$. On $\mathbb{C}^{d-1}\otimes H_{k-1}$ the term $\lambda\Pi_v$ acts as $\lambda h_{k-1}$ on that copy, with
\[
h_{k-1} = b_{k-1}\Bigl(\frac{1}{D_{k-1}} - \frac{1}{D_k}\Bigr) = \frac{b_{k-1}(d-2)}{D_{k-1}(k+d-2)},
\]
and as $\lambda b_{k-1}/D_{k-1}$ on its complement. For $d \ge 3$, $h_{k-1} > 0$. With $X_{\mathrm{op}} = \sum_{k\ge 1} x_k \Pi_k$ and $x_k = g_k^2/(\lambda h_{k-1})$, the operator $I \otimes X_{\mathrm{op}} - \Omega$ consists of the blocks
\[
\begin{pmatrix} x_k & g_k \\ g_k & x_{k-1} + \lambda h_{k-1}\end{pmatrix}, \qquad k = 1, \dots, n,
\]
each tensored with the identity on $\Sym^k(\mathbb{C}^{d-1})$, and is diagonal and non-negative elsewhere. Each block is positive semidefinite, since $x_k \cdot \lambda h_{k-1} = g_k^2$ and $x_{k-1} \ge 0$. Therefore
\[
X - \lambda Y \ \le\ \operatorname{Tr} J_T (I \otimes X_{\mathrm{op}}) = \operatorname{Tr} X_{\mathrm{op}} = \sum_{k=1}^n D_k x_k = \sum_{k=1}^n \frac{k\,b_k}{\lambda(d-2)} = \frac{n s^2}{\lambda(d-2)},
\]
using $D_k/D_{k-1} = (k+d-2)/k$ and $\sum_k k b_k = ns^2$.

\emph{Conclusion.} If $c\gamma \le \varepsilon_\theta/s$ there is nothing to prove. Otherwise $X > 0$, and $\lambda = 2ns^2/(X(d-2))$ gives $X^2(d-2) \le 4ns^2 Y$. With the two tests this is $(d-2)\,s^2(c\gamma - \varepsilon_\theta/s)^2 \le ns^2\varepsilon_\theta$.
\end{proof}

\begin{proof}[Proof of Theorem~\ref{thm:conj}]
Let $\varepsilon \ge \varepsilon_\theta$ be the worst case over the whole family; the right-hand side of Theorem~\ref{thm:conjtheta} decreases in $\varepsilon_\theta$. At $\theta = \pi/4$, $\gamma = a(1-a/2) = (1-e^{-t})/2$ and $c\gamma - \varepsilon/s = (\gamma/2 - \varepsilon)\sqrt{2}$, so $n\varepsilon \ge 2(d-2)(\gamma/2 - \varepsilon)^2 \ge (d-2)\gamma^2/8$ for $\varepsilon \le \gamma/4$. By convexity of the exponential, $1 - e^{-t} \ge (1 - e^{-1})\,t \ge 0.632\,t$ for $t \le 1$, and $1 - e^{-t} \ge 0.632$ for $t \ge 1$; the constants follow.
\end{proof}

\begin{remark}
With $\beta_t = e^{-t/2}(1-e^{-t/2}) \le \gamma$ and $\sin\theta = (\varepsilon/2\beta_t)^{1/3}$, Theorem~\ref{thm:conjtheta} also gives $n\varepsilon \ge (d-2)\beta_t^2(1 - 3(\varepsilon/2\beta_t)^{2/3})_+^2$, whose constant tends to $(d-2)/4$ as $t \to 0$ and $\varepsilon/t \to 0$. The bound is the cost of universal conjugation of the program; compare the optimal fidelity $(n+1)/(n+d)$ of estimating a pure state of dimension $d$ from $n$ copies [3].
\end{remark}

\section{The sharp constant}\label{sec:sharp}

Write $\ell = \ket{L}$ for $L = \proj{0}$, $q_i = \ket{0}\ket{i}\in P$ and $v = \cos\theta\,\ell + e^{i\phi}\sin\theta\, q_i$, so that $\pi_v = \proj{v}$ is the program of $L_v = \ket{0}\bra{\psi}$ with $\psi = \cos\theta\ket{0} + e^{-i\phi}\sin\theta\ket{i}$. Let $\varepsilon$ be the worst-case error of $\Lambda$ over this family, for all $\theta$, $\phi$ and $i = 1, \dots, d-1$.

\begin{theorem}\label{thm:sharpfull}
There is an absolute constant $C$ such that for all $d \ge 2$, $n \ge 1$, $t \in (0, 1/2]$ and every channel $\Lambda$ from $S\otimes P^{\otimes n}$ to $S$,
\[
n\,\varepsilon \ \ge\ \frac{d-1}{4}\, t^2 \Bigl(1 - Ct\bigl(1 + d\,\min\{1, \sqrt{\log(2n/t)/n}\}\bigr)\Bigr).
\]
\end{theorem}

\begin{proof}
\emph{The Choi matrix.} Let $G$ be the map $\rho \otimes \ket{L}\bra{L'} \mapsto L\rho L'^\dagger - \tfrac12(L'^\dagger L\rho + \rho L'^\dagger L)$ on one copy, averaged over the copies with the trace on the others; on $L(S) \otimes L(\Sym^n P)$ it is the first-order part that every protocol must have, since the operators $\pi_v^{\otimes n}$ span $L(\Sym^n P)$. Let $C_0$, $C_1$ be the Choi matrices of $\operatorname{Tr}_P$ and $G$, and $B = \Lambda - \operatorname{Tr}_P - tG$, so the Choi matrix of $\Lambda$ is $C = C_0 + tC_1 + C_B \ge 0$.

\emph{Two test vectors.} Take $e = \ket{0}\ket{\ell}^{\otimes n}\ket{i}$ and $f = -C_0^+C_1 e = (2d\sqrt n)^{-1}\sum_m \ket{m}\ket{\hat\tau}\ket{m}$ with $\hat\tau = n^{-1/2}\sum_j \ket{\ell\cdots q_i\cdots\ell}$ ($q_i$ in copy $j$). Then $e^\dagger C_0 e = e^\dagger C_1 e = 0$, so $e^\dagger C e$ is the population $\bra{i}\Lambda(\proj{0}\otimes\pi_\ell^{\otimes n})\ket{i}$ that leaves $\ket{0}$, which is at most $\varepsilon$ after summing over $i$ since $e^{t\Lind}$ fixes $\proj{0}$. Further $f^\dagger C_0 f = 1/(4n)$, $e^\dagger C_1 f = -1/(4n)$ and $f^\dagger C_1 f = -(d-1+1/n)/(4d^2 n) \le 0$. Positivity of the $2\times 2$ compression gives $e^\dagger C e \ge \lvert e^\dagger C f\rvert^2 / f^\dagger C f$.

\emph{Spin components.} Let $W = \mathrm{span}\{\ell, q_i\}$. The states $\pi_v^{\otimes n}$ lie in $L(\Sym^n W)$, on which $SU(2)$ acting on $W$ decomposes the operators into spins $k = 0,\dots,n$. Write $\pi_v^{\otimes n} = \sum_k O_k(v)$. Each $O_k(v)$ is a rotation of $O_k(\ell)$, with $\lVert O_k\rVert_2^2 = (2k+1)\,n!^2/((n-k)!(n+k+1)!)$, and
\[
\frac{(n+1)\lVert O_k\rVert_2^2}{2k+1} = \prod_{j<k}\frac{n-j}{n+j+2} \le e^{-k(k+1)/(2n+1)} \le e^{-k^2/(2n)} .
\]
Since $O_k$ has rank at most $n+1$, $\tn{O_k(v)} \le \sqrt{2k+1}\,e^{-k^2/(4n)}$ for every $v$. The map $\Lambda(\cdot\otimes O_k(v))$ has diamond norm at most $\tn{O_k}$, and its entries are spherical harmonics of degree $k$ on the Bloch sphere of $W$, whose polar angle is $2\theta$; so they are trigonometric polynomials of degree $2k$ in $\theta$. The Taylor polynomial of order $k_0$ of $e^{t\Lind_v}$ has degree at most $2k_0$ in $\theta$, since $\Lind_v$ is linear in $\proj{v}$.

\emph{Bernstein's inequality.} Let $k_0 = \min\{n+1, \lceil\sqrt{24 n\log(2n/t)}\rceil\}$ and consider $F = \sum_{k \le k_0}\Lambda(\cdot\otimes O_k(v))$ minus the Taylor polynomial of order $k_0$, a trigonometric polynomial of degree at most $2k_0$. It lies within $2\varepsilon + O(t^2/n)$ of zero for all $\theta, \phi$: the discarded spin components and their first two $\theta$-derivatives are bounded by $\sum_{k>k_0}(2k)^r\tn{O_k} \le t^2/n$, $r = 0,1,2$, and the Taylor remainder by $O(t^2(2t)^{k_0}/(k_0+2)!)$. By Bernstein's inequality its first and second $\theta$-derivatives are at most $2k_0$ and $4k_0^2$ times that bound. At $\theta = 0$, the first derivative at two phases and the second derivative averaged over $\phi$ are $B$ applied to $\ket{\hat\tau}\bra{\ell^{\otimes n}}$ (times $\sqrt n$) and to $\proj{\hat\tau}$ (times $2n$, minus $2n\pi_\ell^{\otimes n}$), up to the smooth part of the exponential, which is $O(t^2)$. With $\beta d = 1/(2\sqrt n)$ and $\beta^2 d^2 = 1/(4n)$ for the coefficient $\beta$ of $f$, this gives $\lvert e^\dagger C_B f\rvert \le 2k_0\varepsilon/n + O(t^2/n)$ and $f^\dagger C_B f \le O(k_0^2\varepsilon/n + \varepsilon + t^2)/(4n)$; the factors of $d$ cancel.

\emph{Conclusion.} With $e^\dagger C f = -t/(4n) + e^\dagger C_B f$, the leak towards $\ket{i}$ is at least $\tfrac{t^2}{4n}(1 - O(k_0\varepsilon/t + t))$, where the denominator contributes $O(k_0^2\varepsilon/n + \varepsilon + t^2) \le O(k_0\varepsilon/t + t^2)$ relative to $1/(4n)$. Summing over $i$ bounds $\varepsilon$ from below. If $n\varepsilon > (d-1)t^2/4$ there is nothing to prove; otherwise $k_0\varepsilon/t \le (d-1)t\,k_0/(4n)$, and $k_0/n \le C\min\{1, \sqrt{\log(2n/t)/n}\}$.
\end{proof}

\begin{proof}[Proof of Theorem~\ref{thm:sharp}]
For $t \le c/d$ use $k_0/n \le C$. For the high-accuracy statement, discarding copies shows that the optimal error does not increase with $n$. A protocol with fewer than $N_0 = \lceil A d^2 \log(2d/t)\rceil$ copies therefore has error at least the bound at $N_0$, where $d\sqrt{\log(2N_0/t)/N_0} \le 1$; this exceeds $t^2/(Cd\log(2d/t))$ for a suitable $C$. For $n \ge N_0$ the bound of Theorem~\ref{thm:sharpfull} applies directly.
\end{proof}

\begin{remark}
For $k_0 = n+1$ the argument is the plain Bernstein bound with degree $2n+2$, which gives the statement for $t \le c/d$. The spin decomposition improves on it for $n \gtrsim d^2\log(2n/t)$. On WML's own channel, which uses its $n$ copies in $n$ steps of time $t/n$, the $2\times 2$ bound captures $99.0\%$ to $99.9\%$ of the true leak for $n \le 3$, $d \le 3$ and $t \in [0.0125, 0.05]$. Protocols that use a random number of copies are not covered.
\end{remark}

\section{Constants of wave matrix Lindbladization}\label{sec:wml}

One WML step applies $\exp(\Delta M)$ to $\rho\otimes\pi_L$ for the fixed Lindbladian $M$ of [1, Eqs.~(14), (15)] and traces out the program. By the transfer matrix of [1, App.~B] it equals
\[
\Phi_\Delta = \mathrm{id} + \Delta\Lind + c_J J + c_A A + c_D D_1\circ J,
\]
with $J = L\cdot L^\dagger$, $A = \{L^\dagger L, \cdot\}/2$, $D_1 = (I/d)\operatorname{Tr}$, and
$c_A = \tfrac{2}{d}u^2\varphi(u)$, $c_D = \tfrac{du^2}{d^2-2}\bigl[\tfrac{4(d^2-1)}{d^2}\varphi(v) - 2\varphi(u)\bigr]$, $c_J = -\Delta^2\tfrac{d^2-1}{d}\varphi(v) - \tfrac{2}{d^2}c_D$, where $u = \Delta d/2$, $v = \Delta(d - 1/d)$ and $\varphi(z) = (e^{-z} - 1 + z)/z^2$.

\begin{theorem}[upper bound]\label{thm:wmlupper}
For every $L$ with $\lVert L\rVert_2 = 1$, every $d \ge 2$ and every $\Delta \le 1/(2d)$,
$\tfrac12\dn{e^{\Delta\Lind} - \Phi_\Delta} \le \Delta^2\lVert L\rVert_\infty^2\bigl(d/2 + e^{1/d}\lVert L\rVert_\infty^2\bigr)$. Hence $n \ge \max\{2dt, (d/2 + e^{1/d})t^2/\varepsilon\}$ copies suffice for $\lVert L\rVert_\infty \le 1$.
\end{theorem}

\begin{proof}
The function $\varphi$ is decreasing on $(0,\infty)$ with values in $(0, 1/2]$, and $u \le v \le 1/2$. This gives $0 < c_A, c_D \le d\Delta^2/4$, using $\varphi(v) \ge \varphi(1/2) > 1/3$ for the sign of $c_D$, and $0 < -c_J \le d\Delta^2/2$. The maps $J$ and $D_1\circ J$ are completely positive with adjoint value $L^\dagger L$ at the identity, and $A$ is bounded by $\lVert L^\dagger L\rVert$, so each has diamond norm at most $\kappa = \lVert L\rVert_\infty^2$, and $\dn{\Lind} \le 2\kappa$. The exponential series gives $\dn{e^{\Delta\Lind} - \mathrm{id} - \Delta\Lind} \le e^x - 1 - x \le x^2e^x/2$ with $x = 2\Delta\kappa \le 1/d$. Summing, $\tfrac12\dn{e^{\Delta\Lind} - \Phi_\Delta} \le \tfrac12(d\Delta^2\kappa + 2\Delta^2\kappa^2 e^{1/d})$.
\end{proof}

\begin{theorem}[lower bound for WML]\label{thm:wmllower}
For $L = \proj{0}$ and input $\proj{0}$, the state is fixed by $e^{t\Lind}$, each WML step leaks $\Delta^2(d-1)/4$ in trace distance at leading order, and these errors add up without decay. Hence $n\cdot\mathrm{error}/t^2 \to (d-1)/4$ as $n\to\infty$, for every $t$. For $L = \ket{0}\bra{1}$ and input $\proj{1}$ the one-step error is $\Delta^2(2d-3)/4 + O(\Delta^3)$.
\end{theorem}

So WML needs at least $(\tfrac{d-1}{4} - o(1))\,t^2/\varepsilon$ copies in the worst case as $\varepsilon \to 0$, which is at least $8(d-1)/d \ge 4$ times the constant $d/32$ of [1, Theorem 10], and its worst-case leading constant lies between $d/2 - 3/4$ and $d/2 + 1$. Numerically, the supremum over $L$ of the one-step constant is about $0.300$ at $d = 2$ and $0.750 = (2d-3)/4$ at $d = 3$.

\begin{remark}\label{rem:gap}
For $L = \ket{0}\bra{1}$ and $d \ge 5$, the one-step error $\Delta^2(2d-3)/4$ exceeds $\Delta^2(2d+3)/8$, the bound of [1, Lemma 4], by $7\%$ at $d = 5$ and by a factor approaching $2$ as $d$ grows. This hints at a possible gap in [1, Eq.~(40)]: to leading order $\tilde L - L = (\Delta d/4)A - (\Delta d/2)J$, whose diamond norm for this $L$ is $3\Delta d/4$ rather than $\Delta d/4$. The scaling $O(d\,t^2/\varepsilon)$ of [1] is not affected.
\end{remark}

\section{Checks and limits}

The scripts in \path{laif/wml-sample-constants/sim} check, with asserts, the program convention and the coherence bound of Theorem~\ref{thm:conjtheta} for $t$ from $0.01$ to $12$, positivity and trace of its certificate from exact Haar moments for $(d,n)$ up to $(3,4)$, $(4,3)$, $(5,2)$ and $(6,2)$, the spin decomposition of Theorem~\ref{thm:sharpfull} ($\lVert O_k\rVert_2$ exactly for $n \le 100$, the tails beyond $k_0$ for $n \le 2000$, and the degree of each component under a random channel), the identities and the $2\times 2$ bound behind its first form, and every coefficient identity of WML against the matrix exponential. The constants of Theorem~\ref{thm:sharpfull} are not optimized. Theorem~\ref{thm:conj} has the factor $d-2$ and says nothing at $d = 2$, where Theorem~\ref{thm:sharp} applies; its constant $0.0124$ is far below $(d-1)/4$. Between the accuracies $t^2/(Cd\log(2d/t))$ and $(1-e^{-t})/8$ for $c/d < t \le c$, the factor $d$ follows from Theorem~\ref{thm:conj} but the constant $(d-1)/4$ is open.

\section*{Use of AI}
The results, proofs, scripts and Lean files were developed with the help of an AI system (Claude Opus 5.5), which also checked them. The author takes responsibility for the content.

\section*{Code}
Scripts and Lean files: \url{https://github.com/ssenge/LAIF}, folder \path{laif/wml-sample-constants}.

\section*{References}
\begin{enumerate}
\item[{[1]}] S.~Park, Y.~Seo, B.~Go, D.~Patel, M.~M.~Wilde, H.~Kwon. Improved sample complexity bound for sample-based Lindbladian simulation. arXiv:2605.30301 (2026).
\item[{[2]}] B.~Go, H.~Kwon, S.~Park, D.~Patel, M.~M.~Wilde. Sample-based Hamiltonian and Lindbladian simulation: non-asymptotic analysis of sample complexity. arXiv:2412.02134.
\item[{[3]}] D.~Bru{\ss}, C.~Macchiavello. Optimal state estimation for $d$-dimensional quantum systems. Phys.\ Lett.\ A 253, 249 (1999).
\item[{[4]}] Wave matrix Lindbladization II: general Lindbladians, linear combinations, and polynomials. arXiv:2309.14453 (2023).
\end{enumerate}

\end{document}
